Une tige uniforme de \qty{\v{M}}{kg} et de \qty{\v{L}}{m} peut pivoter sans frottement autour d'un axe horizontal passant par une de ses extrémités. On la tient à l'horizontale, puis on la lâche du repos ($I_{\text{bout}} = \tfrac{1}{3}ML^2$).

\begin{center}
\begin{tikzpicture}[thick]
    \fill[pattern=north east lines] (-0.55,-1.4) rectangle (-0.3,0.8);
    \draw (-0.3,-1.4) -- (-0.3,0.8);
    \draw (-0.3,0) -- (0,0);
    \draw[fill=gray!50] (0,-0.08) rectangle (5,0.08);
    \draw[fill=white] (0,0) circle (2.5pt);
    \draw[thin,|-|] (0,0.45) -- node[above] {\small\qty{\v{L}}{m}} (5,0.45);
    \draw[thin,-{Stealth[length=2.5mm]}] (0,0) ++(-4:4.6) arc (-4:-30:4.6);
\end{tikzpicture}
\end{center}

\begin{enumerate}[label=\alph*)]
    \item Déterminez le moment de force exercé par la gravité au moment où on lâche la tige.
    \item Déterminez l'accélération angulaire initiale de la tige.
    \item Déterminez la vitesse de l'extrémité libre lorsque la tige passe à la verticale.
\end{enumerate}
